\section{The finite flat critical divisor and its ramification}
\label{sec:projective-critical-v144}

The critical cover of Section~\ref{sec:critical-correspondence-v143}
was written on an affine spectral chart and away from its discriminant.
We now retain the full projective critical divisor.  It remains finite
and flat when critical directions cross the omitted chart or collide.
Its nonreduced fibres and its ramification are recovered from the same
finite failure invariant.  Thus the spectral construction carries more
than the generic critical count or its real realization.

\subsection{The projective score differential}
Fix positive integers \(m_1,\ldots,m_r\), with \(r\ge2\) and
\(n=\sum_i m_i\).  On the split semisimple locus write
\(K=xB+yA\), where \(A\) is nonsingular and \(A^{-1}B\) has
distinct spectral values \(\alpha_i\) with these multiplicities.
Over a nonreduced base, this means that mutually orthogonal idempotent
projectors \(P_i\) of ranks \(m_i\) are supplied, with
\(\sum_iP_i=1\) and \(A^{-1}B=\sum_i\alpha_iP_i\).
It is stronger than semisimplicity of the geometric fibres alone.
The base field is \(\mathbb C\); all the universal formulas below
are already defined over \(\mathbb Q\).  Put
\begin{equation}\label{eq:binary-spectral-data-v144}
 \begin{aligned}
 L_i&=\alpha_i x+y,&
 D_h&=\prod_i L_i,& D_i&=\prod_{j\ne i}L_j,\\
 Q_h&=\sum_i\sigma_iD_i,&
 U(K)&=\operatorname{tr}(SK^{-1})=Q_h/D_h.
 \end{aligned}
\end{equation}
Here \(\sigma_i=\operatorname{tr}(P_iA^{-1}S)\), with the
projectors from \eqref{eq:spectral-projectors-v143}.  The polynomial
\(Q_h\) is a binary form of degree \(r-1\); its affine leading
coefficient is allowed to vanish.

Let \(\mathscr U\) be the open subscheme of
\(\Spec\mathbb Q[\alpha_1,\ldots,\alpha_r,\sigma_1,\ldots,\sigma_r]\)
on which the \(\alpha_i\) are pairwise distinct, \(Q_h\) has
\(r-1\) distinct projective zeros, and these zeros are disjoint from
those of \(D_h\).  Equivalently, invert the spectral differences,
the binary discriminant of \(Q_h\), and
\(\operatorname{Res}(D_h,Q_h)\).  The discriminant of a binary
linear form is understood as \(1\); the resultant condition then
also excludes the zero form.  No discriminant of a score polynomial
is inverted.  In particular, this base allows critical-point
collisions.  All subsequent base changes are from \(\mathscr U\).

Consider the following rational one-form on the projective line:
\begin{equation}\label{eq:projective-score-form-v144}
 \begin{aligned}
 \omega&=\sum_i(n-m_i)\,d\log L_i-n\,d\log Q_h
       =\frac{\eta}{D_hQ_h},\\
 \eta&=Q_h\sum_i(n-m_i)D_i\,dL_i-nD_h\,dQ_h.
 \end{aligned}
\end{equation}
Differentials are relative to the base.  The equality
\(\sum_i(n-m_i)=n(r-1)\) makes this expression independent
of rescaling a representative \((x,y)\).  If
\(\eta=\eta_x\,dx+\eta_y\,dy\), Euler's identity gives
\(x\eta_x+y\eta_y=0\).  The syzygy of \((x,y)\) therefore
gives a unique binary form \(G\) of degree \(2r-3\) such that
\begin{equation}\label{eq:binary-score-v144}
                   \eta=G(x,y)(x\,dy-y\,dx).
\end{equation}
This argument is valid over the universal coefficient ring: reducing
\(x\eta_x+y\eta_y=0\) modulo \(x\) first shows that
\(x\) divides \(\eta_y\), and then determines both coefficients.
It remains valid after every base change.

On the chart \((x,y)=(1,-t)\), let
\(D(t)=\prod_i(t-\alpha_i)\) and
\(Q(t)=\sum_i\sigma_iD(t)/(t-\alpha_i)\).  Direct substitution
in \eqref{eq:projective-score-form-v144} gives
\begin{equation}\label{eq:projective-affine-score-v144}
 G(1,-t)=-F(t),\qquad
 F=nDQ'-Q\sum_i(n-m_i)\frac{D}{t-\alpha_i}.
\end{equation}
Thus the form \(G\), rather than the degree of a particular affine
polynomial, records the whole critical divisor.

\begin{theorem}[Finite flat reconstruction of the critical divisor]
\label{thm:projective-critical-v144}
Over \(\mathscr U\), the zero scheme
\(\mathscr Z=V(G)\subset\mathbb P^1_{\mathscr U}\) is finite
locally free of degree \(2r-3\), and is disjoint from
\(V(D_hQ_h)\).  It is isomorphic, as a scheme over the base, to
the entire reciprocal critical scheme on invertible pencil members.
These assertions include nonreduced fibres and arbitrary base changes.

Under a change of pencil frame the divisor transforms on the
projective pencil line; under simultaneous congruence it is unchanged
up to the induced isomorphism.  For a family specified by its framed
first relation spaces, with split semisimple spectral data and data
matrices as above, the divisor is consequently reconstructed from
its order-\(d\) failure neighbourhoods.  Here
\(d=n^2+2n-4\), and the supplied data and semisimple condition are
additional to, not restrictions on, the all-pencil inverse theorem.
\end{theorem}
\begin{proof}
At a zero of \(L_i\) the residue of \(\omega\) is \(n-m_i\);
at a zero of \(Q_h\) it is \(-n\).  These are nonzero units in
characteristic zero.  All the poles are simple and mutually disjoint
on \(\mathscr U\).  Hence in every geometric fibre the numerator
\(G\) is nonzero at every pole.  This proves disjointness as schemes:
a nonempty intersection would have a geometric point in some fibre.
It also proves that \(G\) is not identically zero in any fibre.

Here is an elementary finite-flat argument which also deals with
arbitrary, possibly nonreduced, bases.  At each point of
\(\mathscr U\), choose a constant rational point of the projective
line outside the finite zero set of its fibre of \(G\).  The
characteristic-zero residue field contains its prime field \(\mathbb Q\).
Thus \(\mathbb P^1(\mathbb Q)\) supplies infinitely many constant
points, whereas the nonzero fibre form has only finitely many zeros.
One of these points avoids them.  Its evaluation is nonzero in the
residue field, and becomes a unit after shrinking the base.
A constant projective coordinate change takes this point to infinity.
The scheme \(\mathscr Z\) is then contained in the complementary
affine chart, where its equation has unit leading coefficient and
can be made monic of degree \(2r-3\).  Its algebra is free on
\(1,t,\ldots,t^{2r-4}\).  These neighbourhoods cover the base,
proving finite local freeness and the stated degree.  The same
presentations survive all base changes.  They also identify
\(\mathscr Z\) as a relative effective Cartier divisor; compare
the general criterion in \cite[Tag~062Y]{StacksCritical144}.

We next prove the identification with the critical scheme, rather
than only a bijection of closed points.  Locally trivialize the
projective direction by a nonsingular representative \(K_0(t)\)
and write \(K=\lambda K_0(t)\), where \(\lambda\) is invertible.
Let \(U_0(t)=\operatorname{tr}(S K_0(t)^{-1})\).  The rational
differential of the reciprocal likelihood is obtained from
\[
 -n\log\lambda-\log\det K_0(t)-U_0(t)/\lambda.
\]
Its radial score is \((-n\lambda+U_0)/\lambda^2\).  Thus the
critical coordinate ring imposes exactly
\(\lambda=U_0/n\), and makes \(U_0\) a unit.  Eliminating
\(\lambda\) from the remaining score gives
\[
             -d\log\det K_0-n\,d\log U_0=\omega.
\]
The determinant factorization and \(U_0=Q_h/D_h\) give the
last equality.  Since \(\mathscr Z\) is disjoint from \(D_hQ_h\),
all divisions are by units on \(\mathscr Z\).  These are inverse
coordinate-ring homomorphisms, including nilpotents.  Globally the
inverse sends a direction represented by \(K_0\) to
\begin{equation}\label{eq:radial-reconstruction-v144}
                         K=\frac{U(K_0)}{n}K_0.
\end{equation}
Replacing \(K_0\) by \(aK_0\) replaces \(U(K_0)\) by
\(a^{-1}U(K_0)\), so this formula glues.  It also proves that no
projective chart has been omitted.

Both \(-d\log\det K_0-n\,d\log U(K_0)\) and
\eqref{eq:radial-reconstruction-v144} are invariant under pencil
reparametrization.  Simultaneous congruence preserves the trace and
changes the determinant by a constant square; its relative logarithmic
differential is therefore unchanged.  Finally apply
Theorem~\ref{thm:finite-parameter-immersion} and the universal
construction in Theorem~\ref{thm:critical-correspondence-v143}.
This last reconstruction step uses a framed classifying family;
it asserts neither an equivalence of unmarked moduli stacks nor
recovery of a data matrix from an unmarked complex scheme alone.
\end{proof}

\subsection{Hessian degeneracy as scheme-theoretic ramification}
Write \(\pi:\mathscr Z\to\mathscr U\).  The relative ramification
scheme is defined by
\(\operatorname{Fitt}_0\Omega_{\mathscr Z/\mathscr U}\).
The trace discriminant on the base is the determinant of
\((a,b)\mapsto\operatorname{Tr}_{\pi}(ab)\) on
\(\pi_*\mathcal O_{\mathscr Z}\); it is an ideal locally generated
by a section of an invertible sheaf, not a choice of eigenbasis.

\begin{theorem}[Ramification, Hessians, and the trace discriminant]
\label{thm:critical-ramification-v144}
Under Theorem~\ref{thm:projective-critical-v144}, the ramification
ideal is exactly the ideal generated by the determinant of the
reciprocal score Hessian on the critical scheme.  The trace-discriminant
ideal is locally generated, up to a unit, by the binary discriminant
\(\operatorname{Disc}(G)\).  Both ideals commute with arbitrary base
change.  The finite critical scheme is etale precisely off this
trace discriminant, and the totally real cover of
Theorem~\ref{thm:real-critical-cover-v143} is a restriction of this
finite flat family.
\end{theorem}
\begin{proof}
On a base neighbourhood with the monic presentation
\(\mathcal A=\mathcal O_{\mathscr U}[t]/(g(t))\) used above,
\[
       \Omega_{\mathcal A/\mathcal O_{\mathscr U}}
             =\mathcal A\,dt/(g'(t)\,dt).
\]
Its zeroth Fitting ideal is \((g')\).  In the local coordinates
\((\lambda,t)\) of the preceding proof the radial second derivative,
after imposing the radial score, is \(-n/\lambda^2\), a unit.
The complementary Schur factor of the Hessian is the derivative
of the eliminated score.  That score is \(v(t)g(t)\), with
\(v\) a unit along the critical scheme, since its only denominators
come from the disjoint poles.  Its derivative modulo \((g)\) is
\(v g'\).  The Hessian determinant thus generates \((g')\).
A change of smooth pencil coordinates multiplies this determinant
on the critical scheme by the square of a Jacobian unit.  This
proves the claimed equality of ideal sheaves independently of the
coordinates and without reducing the scheme.

In the basis \(1,t,\ldots,t^{m-1}\), \(m=2r-3\), the
determinant of the trace pairing for a monic polynomial algebra
is \(\operatorname{disc}(g)\).  One proof checks the Vandermonde
identity for the polynomial with independent roots and then descends
the resulting polynomial identity in its coefficients; consequently
it remains valid for repeated roots and arbitrary coefficient rings.
Making the leading coefficient a unit and changing projective frame
multiply the binary discriminant by units.  Hence these local ideals
are exactly \((\operatorname{Disc}(G))\).  This uses the classical
trace-discriminant construction, not a new general discriminant
theorem; see \cite[Tags~0BVH, 0BJF]{StacksCritical144}.

The differential presentation, its Fitting ideal, and the matrix of
the trace pairing all commute with base change.  A monic algebra is
etale exactly when \(g'\) is a unit in it, equivalently when its
discriminant is a unit.  On the separated real chamber the sign
argument in Theorem~\ref{thm:real-critical-cover-v143} proves this
condition and labels all \(m\) real roots.  Its additional affine
leading-coefficient localization merely selects one chart of the
present family.  No new sign argument is needed.
\end{proof}

The trace discriminant on the base is not being identified with an
arbitrarily scheme-theoretically imaged ramification locus: their
multiplicities need not coincide.  Likewise finite flatness is asserted
on \(\mathscr U\), not over data for which a zero of \(Q_h\)
collides with a pole or becomes multiple.  These exclusions concern
the spectral data polynomial, not collisions of zeros of \(G\).

\paragraph{Coordinate invariance on the critical quotient.}
Write a smooth local coordinate change as \(x=x(y)\), its Jacobian as
\(J=(\partial x_a/\partial y_i)\), and the first scores as \(s_a\).
The second-derivative chain rule is
\[
 H_y=J^{\mathsf t}H_xJ+
             \sum_a s_a\left(\frac{\partial^2x_a}
                                   {\partial y_i\partial y_j}\right)_{ij}.
\]
All \(s_a\) vanish in the critical quotient itself, not only on its
reduction.  There \(H_y=J^{\mathsf t}H_xJ\) and
\(\det H_y=(\det J)^2\det H_x\).  Since \(\det J\) is a unit,
the Hessian ideal is unchanged even in a nonreduced critical algebra.

\subsection{A length-three critical fibre at a nonsingular pencil member}
\begin{proposition}[A triple critical collision in a fixed pencil]
\label{prop:triple-critical-v144}
Let \(A=I_3\), \(B=\operatorname{diag}(-1,0,1)\), and
\[
 S_c=\operatorname{diag}\bigl((1+c)/2,-c,(1+c)/2\bigr),
       \qquad c\in\mathbb C\setminus\{0,-1\}.
\]
The reciprocal critical schemes form a finite flat family of length
three.  At \(c=1\) all three critical points are real and reduced;
one lies in the projective direction \(x=0\).  At \(c=-2/3\)
the full critical scheme is supported at \(K=I_3/3\) and is
isomorphic to \(\Spec\mathbb C[u]/(u^3)\).  The trace-discriminant
ideal of this family is generated, up to a unit, by
\((3c+2)^3\).
\end{proposition}
\begin{proof}
Here \(D(t)=t(t^2-1)\), \(Q(t)=t^2+c\), and
\begin{equation}\label{eq:triple-score-v144}
                F(t)=-(6c+4)t^2+2c.
\end{equation}
The excluded parameters are exactly those where \(Q\) is repeated
or meets a zero of \(D\).  Thus this entire family lies in
\(\mathscr U\).  Homogenizing \(F\) to its projective degree
three gives
\(\widehat F(T,Z)=-(6c+4)T^2Z+2cZ^3\).
The associated projective direction is \([x:y]=[Z:-T]\).
At \(c=1\) its zeros are \(t=\pm1/\sqrt5\) and the point
\([T:Z]=[1:0]\), all simple.  Formula
\eqref{eq:radial-reconstruction-v144} takes the latter to \(I_3/3\).

At \(c=-2/3\), \(\widehat F=-4Z^3/3\), so its whole
zero scheme is in the chart \(T=1\), where \(u=Z/T=-x/y\).
The local equation is \(-4u^3/3\).  The radial reconstruction is
an isomorphism on this nonreduced scheme and its reduced point is
\(I_3/3\), since \(\operatorname{tr}S_c=1\).
The binary cubic discriminant, or the trace calculation in a monic
chart, is
\begin{equation}\label{eq:triple-discriminant-v144}
                 \operatorname{Disc}(\widehat F)
                       =64c(3c+2)^3.
\end{equation}
Since \(c\) is a unit on the stated base, this gives the claimed
ideal.  The special ramification ideal is \((u^2)\) in
\(\mathbb C[u]/(u^3)\), in agreement with the Hessian theorem.
\end{proof}

This example keeps the pencil and its spectral sheaf fixed.  It shows
why the critical correspondence must be retained with its scheme
structure: the number of distinct critical points drops while the
finite flat degree stays three.  The order-\(d\) inverse theorem,
followed by the supplied data family, reconstructs this collision and
its trace-discriminant multiplicity.  This is a structural consequence
of the same reconstruction chain, not a replacement for the independent
real-root argument or a claim that generic degree formulas are new.
